\paragraph{Opened HLP source and its precise scope.}
The copy actually opened is the original 1934 edition of
\emph{Inequalities}, not the unverified 1952 second edition.
Chapter~X, Theorems~368--370 (\S\S10.2--10.4, pp.~261--265),
gives the finite scalar-product rearrangement and two-factor
coefficient comparisons corresponding to the rearrangement part of
Lemma~\ref{C:majorization}.  For its additional convolution comparison,
write \(a^+\) for the values of a finitely supported nonnegative
sequence on \(\mathbb Z\), placed in decreasing order at
\(0,1,-1,2,-2,\ldots\), and put \(a^-(j)=a^+(-j)\).
The source calls a nonzero multiset symmetric when its largest
positive value has odd multiplicity and its other positive values
have even multiplicities; zeros may be added to obtain the book's
odd-length indexing interval.  In that case \(a^+=a^-=a^\star\)
is even and decreases with \(|j|\).
With \(a,b\) arbitrary and each of \(c,d,\ldots\) symmetric in
this sense, Theorem~373 and its extension, Theorem~374, state
\[
 (a*b*c*d*\cdots)(0)
 \ \leq\
 (a^+*b^-*c^\star*d^\star*\cdots)(0).
\]
The complete finite proof is in \S\S10.5--10.9, pp.~265--274:
it reduces to indicator layers, proves the three-factor comparison,
and uses closure of symmetric decreasing convolution
(Theorem~375, p.~273) for the induction.
These exact source hypotheses are retained here
\cite[\S\S10.1--10.9, pp.~260--274]{C_HLP1934}.

This source comparison concerns rearrangements of fixed finite
multisets.  The separate positivity and convex reduction in
Lemma~\ref{C:majorization} is the justification used in this
manuscript for passing from suitable rearrangement inequalities to
weak majorization; one cannot instead assume that an arbitrary
weakly majorized input retains the displayed symmetry.
Nor can the integer coefficient relation in HLP be transferred
to an arbitrary group: Theorem~\ref{C:K-counterexample} gives an
explicit obstruction to that proposed Condition~K implication.
HLP \S10.10 itself retains only a multiplicative constant depending
on the number of factors in its unrestricted Theorem~377
(pp.~274--276), rather than the same prescribed-order
constant-one comparison.

Finally, these inequalities do not construct the labelled chains,
the map \(\rho\), its specified kernel, or the covariance identities
of the manuscript's system.  Proposition~\ref{C:integer-system}
checks those separately, including intervals of even cardinality,
which have half-integer centres.  An even-cardinality indicator
does not satisfy the symmetric-multiset hypothesis above, and
Theorem~374 allows only two arbitrary multisets.  Thus these quoted
HLP results alone do not establish Condition~K for the full integer
system with arbitrarily many such inputs.  This is a limitation of
the cited comparison, not an assertion that Condition~K fails for
that integer system.  The original Gabriel article has now also
been opened: Theorems~1--2 (pp.~32--39) give the finite three-factor
and additional-symmetric-factor comparisons above, while
Theorems~3--4 (pp.~39--42) concern paired autocorrelations and
products of squared moduli of finite trigonometric polynomials
\cite{C_Gabriel1932}. Their paired structure does not supply a
comparison for unrestricted independent factors, and none of these
results supplies the manuscript's additional system axioms.
